% GENERATED FILE. Edit canonical lessons, then run npm run generate:books.
% canonical-source-hash: fnv1a64:8489d87525e2a3f8
% canonical-lessons: ML-C535-arbudam, ML-C535-mutantu, ML-C535-cennikkuttu, ML-C535-vatam, ML-C535-charddikkuka

\chapter{Aches That Linger}
\label{ch:ml-aches-that-linger}
\hlchaptermodality{535}

\begin{chapteropening}
\textbf{By the end of this chapter:} \emph{I can say cancer, a limp, a migraine, rheumatism and to throw up.}

Complete the last lesson of chapter 535: I can say cancer, a limp, a migraine, rheumatism and to throw up.
\end{chapteropening}

\section[\texorpdfstring{\ml{അർബുദം} (arbudaṁ)}{അർബുദം (arbudaṁ)}]{\ml{അർബുദം} (arbudaṁ) — cancer}

\label{lesson:ML-C535-arbudam}



\begin{quote}
Before the new one: say the Malayalam for insomnia, then the Malayalam for to have a cough.
\end{quote}

\subsection*{You'll want to know: \ml{അർബുദം}}
\textbf{\ml{അർബുദം}} — \emph{arbudaṁ} — \textquotedblleft{}cancer\textquotedblright{}.

\begin{quote}
\ml{പുകവലി} \ml{അർബുദത്തിന്} \ml{കാരണമാകും}.
\end{quote}

\begin{quote}
\emph{pukavali arbudattinŭ kāraṇamākuṁ.}
\end{quote}

\begin{quote}
Smoking can cause cancer.
\end{quote}

In speech many say \ml{ക്യാൻസർ}.

\begin{cousinweb}
\textbf{\ml{അർബുദം}} (\emph{arbudaṁ}) is borrowed from Sanskrit \emph{arbuda}, \textquotedblleft{}a long round mass; a tumour\textquotedblright{} (CDIAL 686).
\end{cousinweb}

\begin{grammarlens}[title={What you've built}]
One more word for describing what you lived and felt.
\end{grammarlens}

\subsection*{Your turn}
\begin{itemize}
\raggedright
  \item \emph{Say it:} \emph{arbudaṁ}
  \item \emph{Say it:} \emph{arbudaṁ}, once more
  \item \emph{Say it:} \emph{pukavali arbudattinŭ kāraṇamākuṁ.} — Smoking can cause cancer.
  \item \emph{Recall:} say the Malayalam for insomnia, then the Malayalam for to have a cough, then say \emph{arbudaṁ} again
\end{itemize}

\begin{tcolorbox}[breakable,colback=teal!4,colframe=teal!35!black,title={Before you move on}]
What does \ml{അർബുദം} mean? (\textquotedblleft{}Cancer\textquotedblright{}.) Say it once more.
\end{tcolorbox}
\section[\texorpdfstring{\ml{മുടന്ത്} (muṭantŭ)}{മുടന്ത് (muṭantŭ)}]{\ml{മുടന്ത്} (muṭantŭ) — a limp; lameness}

\label{lesson:ML-C535-mutantu}



\begin{quote}
Before the new one: say the Malayalam for to have a cough, then the Malayalam for cancer.
\end{quote}

\subsection*{You'll want to know: \ml{മുടന്ത്}}
\textbf{\ml{മുടന്ത്}} — \emph{muṭantŭ} — \textquotedblleft{}a limp; lameness\textquotedblright{}.

\begin{quote}
\ml{കാലിലെ} \ml{മുറിവ്} \ml{കാരണം} \ml{അവന്} \ml{മുടന്ത്} \ml{ഉണ്ട്}.
\end{quote}

\begin{quote}
\emph{kālile muṟivŭ kāraṇaṁ avanŭ muṭantŭ uṇṭŭ.}
\end{quote}

\begin{quote}
He has a limp because of the wound on his leg.
\end{quote}

\ml{മുടന്തുക} is to limp.

\begin{cousinweb}
\textbf{\ml{മുടന്ത്}} (\emph{muṭantŭ}) is native Dravidian: DEDR 4919 lists Malayalam words for lameness and limping, among them \ml{മുടന്തൻ}, a lame man (DEDR 4919).
\end{cousinweb}

\begin{grammarlens}[title={What you've built}]
One more word for describing what you lived and felt.
\end{grammarlens}

\subsection*{Your turn}
\begin{itemize}
\raggedright
  \item \emph{Say it:} \emph{muṭantŭ}
  \item \emph{Say it:} \emph{muṭantŭ}, once more
  \item \emph{Say it:} \emph{kālile muṟivŭ kāraṇaṁ avanŭ muṭantŭ uṇṭŭ.} — He has a limp because of the wound on his leg.
  \item \emph{Recall:} say the Malayalam for to have a cough, then the Malayalam for cancer, then say \emph{muṭantŭ} again
\end{itemize}

\begin{tcolorbox}[breakable,colback=teal!4,colframe=teal!35!black,title={Before you move on}]
What does \ml{മുടന്ത്} mean? (\textquotedblleft{}A limp; lameness\textquotedblright{}.) Say it once more.
\end{tcolorbox}
\section[\texorpdfstring{\ml{ചെന്നിക്കുത്ത്} (cennikkuttŭ)}{ചെന്നിക്കുത്ത് (cennikkuttŭ)}]{\ml{ചെന്നിക്കുത്ത്} (cennikkuttŭ) — a migraine}

\label{lesson:ML-C535-cennikkuttu}



\begin{quote}
Before the new one: say the Malayalam for cancer, then the Malayalam for a limp.
\end{quote}

\subsection*{You'll want to know: \ml{ചെന്നിക്കുത്ത്}}
\textbf{\ml{ചെന്നിക്കുത്ത്}} — \emph{cennikkuttŭ} — \textquotedblleft{}a migraine\textquotedblright{}.

\begin{quote}
\ml{വെയിലത്ത്} \ml{നടന്നാൽ} \ml{അമ്മയ്ക്ക്} \ml{ചെന്നിക്കുത്ത്} \ml{വരും}.
\end{quote}

\begin{quote}
\emph{veyilattŭ naṭannāl ammaykkŭ cennikkuttŭ varuṁ.}
\end{quote}

\begin{quote}
If she walks in the sun, Mother gets a migraine.
\end{quote}

\ml{ചെന്നി} the temple of the head + \ml{കുത്ത്} a stab: the pain that stabs at one side.

\begin{cousinweb}
\textbf{\ml{ചെന്നിക്കുത്ത്}} (\emph{cennikkuttŭ}) is native Dravidian, a compound of \ml{ചെന്നി}, \textquotedblleft{}the temples\textquotedblright{} (DEDR 1989), and \ml{കുത്ത്}, \textquotedblleft{}a stab, a piercing pain\textquotedblright{} (DEDR 1719).
\end{cousinweb}

\begin{grammarlens}[title={What you've built}]
One more word for describing what you lived and felt.
\end{grammarlens}

\subsection*{Your turn}
\begin{itemize}
\raggedright
  \item \emph{Say it:} \emph{cennikkuttŭ}
  \item \emph{Say it:} \emph{cennikkuttŭ}, once more
  \item \emph{Say it:} \emph{veyilattŭ naṭannāl ammaykkŭ cennikkuttŭ varuṁ.} — If she walks in the sun, Mother gets a migraine.
  \item \emph{Recall:} say the Malayalam for cancer, then the Malayalam for a limp, then say \emph{cennikkuttŭ} again
\end{itemize}

\begin{tcolorbox}[breakable,colback=teal!4,colframe=teal!35!black,title={Before you move on}]
What does \ml{ചെന്നിക്കുത്ത്} mean? (\textquotedblleft{}A migraine\textquotedblright{}.) Say it once more.
\end{tcolorbox}
\section[\texorpdfstring{\ml{വാതം} (vātaṁ)}{വാതം (vātaṁ)}]{\ml{വാതം} (vātaṁ) — rheumatism; joint pain}

\label{lesson:ML-C535-vatam}



\begin{quote}
Before the new one: say the Malayalam for a limp, then the Malayalam for a migraine.
\end{quote}

\subsection*{You'll want to know: \ml{വാതം}}
\textbf{\ml{വാതം}} — \emph{vātaṁ} — \textquotedblleft{}rheumatism; joint pain\textquotedblright{}.

\begin{quote}
\ml{മഴക്കാലത്ത്} \ml{മുത്തച്ഛന്} \ml{വാതം} \ml{കൂടും}.
\end{quote}

\begin{quote}
\emph{maḻakkālattŭ muttacchanŭ vātaṁ kūṭuṁ.}
\end{quote}

\begin{quote}
In the rainy season grandfather's rheumatism gets worse.
\end{quote}

In traditional medicine \ml{വാതം} is also one of the three humours of the body.

\begin{cousinweb}
\textbf{\ml{വാതം}} (\emph{vātaṁ}) is borrowed from Sanskrit \emph{vāta}, \textquotedblleft{}wind\textquotedblright{}, and in medicine \textquotedblleft{}wind as a humour of the body\textquotedblright{} (CDIAL 11491).
\end{cousinweb}

\begin{grammarlens}[title={What you've built}]
One more word for describing what you lived and felt.
\end{grammarlens}

\subsection*{Your turn}
\begin{itemize}
\raggedright
  \item \emph{Say it:} \emph{vātaṁ}
  \item \emph{Say it:} \emph{vātaṁ}, once more
  \item \emph{Say it:} \emph{maḻakkālattŭ muttacchanŭ vātaṁ kūṭuṁ.} — In the rainy season grandfather's rheumatism gets worse.
  \item \emph{Recall:} say the Malayalam for a limp, then the Malayalam for a migraine, then say \emph{vātaṁ} again
\end{itemize}

\begin{tcolorbox}[breakable,colback=teal!4,colframe=teal!35!black,title={Before you move on}]
What does \ml{വാതം} mean? (\textquotedblleft{}Rheumatism; joint pain\textquotedblright{}.) Say it once more.
\end{tcolorbox}
\section[\texorpdfstring{\ml{ഛർദ്ദിക്കുക} (charddikkuka)}{ഛർദ്ദിക്കുക (charddikkuka)}]{\ml{ഛർദ്ദിക്കുക} (charddikkuka) — to throw up; to vomit}

\label{lesson:ML-C535-charddikkuka}



\begin{quote}
Before the new one: say the Malayalam for a migraine, then the Malayalam for rheumatism.
\end{quote}

\subsection*{You'll want to know: \ml{ഛർദ്ദിക്കുക}}
\textbf{\ml{ഛർദ്ദിക്കുക}} — \emph{charddikkuka} — \textquotedblleft{}to throw up; to vomit\textquotedblright{}.

\begin{quote}
\ml{ബസ്} \ml{യാത്രയിൽ} \ml{കുട്ടി} \ml{ഛർദ്ദിച്ചു}.
\end{quote}

\begin{quote}
\emph{basŭ yātrayil kuṭṭi charddiccu.}
\end{quote}

\begin{quote}
The child threw up on the bus journey.
\end{quote}

Past: \ml{ഛർദ്ദിച്ചു}, \emph{charddiccu}.

The noun \ml{ഛർദ്ദി}, vomiting, is already familiar.

\begin{cousinweb}
\textbf{\ml{ഛർദ്ദിക്കുക}} (\emph{charddikkuka}) is built from Sanskrit \emph{chardi}, \textquotedblleft{}vomiting\textquotedblright{} (CDIAL 4999), with the Malayalam verb ending -\ml{ിക്കുക}.
\end{cousinweb}

\begin{grammarlens}[title={What you've built}]
One more word for describing what you lived and felt.
\end{grammarlens}

\subsection*{Your turn}
\begin{itemize}
\raggedright
  \item \emph{Say it:} \emph{charddikkuka}
  \item \emph{Say it:} \emph{charddikkuka}, once more
  \item \emph{Say it:} \emph{basŭ yātrayil kuṭṭi charddiccu.} — The child threw up on the bus journey.
  \item \emph{Recall:} say the Malayalam for a migraine, then the Malayalam for rheumatism, then say \emph{charddikkuka} again
\end{itemize}

\begin{tcolorbox}[breakable,colback=teal!4,colframe=teal!35!black,title={Before you move on}]
What does \ml{ഛർദ്ദിക്കുക} mean? (\textquotedblleft{}To throw up; to vomit\textquotedblright{}.) Say it once more.
\end{tcolorbox}
